\section{导读：AI 社交不是“AI 陪你聊天”}

这期访谈围绕自然选择的两款产品：AI 游戏 Eve 和 AI 社交 Elys。Tristan 给出的核心叙事是“人类太孤独了”，但这并不是单纯的情绪陪伴产品。Elys 的更大野心，是让 AI 分身主动面对世界，帮助用户把自己的 context 转化为真实连接。它试图回答一个互联网时代遗留问题：社交网络为什么总是把复杂的人压缩成低维标签，而 AI 能否第一次让高维 context 流动起来。

\begin{importantbox}{本期核心命题}
传统互联网社交依赖低维标签和用户主动行为；AI 社交网络试图让用户构建主体性，让赛博分身带着 context 主动表达、推荐、破冰，并最终带回真人连接。
\end{importantbox}

\begin{knowledgebox}{视觉策略说明}
本视频是固定访谈画面，没有教学 slides、白板或产品演示。按本仓库播客标准，正文不重复插入人物帧；封面用于来源识别，正文用概念图和表格承载产品逻辑。
\end{knowledgebox}

\subsection{本章小结}

Elys 不是一个简单的 AI Tinder，也不是 AI 和 AI 互相聊天的玩具。它的关键变量是 context、主体性、隐私交换、真人连接率和 proactive 产品形态。

